%========================================================
% PTUK GRADUATION PROJECT REPORT
% Project: Ardi - أرضي
%========================================================

\documentclass[12pt,a4paper,oneside]{report}

%========================================================
% PACKAGES
%========================================================

\usepackage{fontspec}
\usepackage{polyglossia}

\usepackage[
    a4paper,
    top=2.5cm,
    bottom=2.5cm,
    left=3cm,
    right=2.5cm,
    headheight=18pt,
    headsep=0.6cm,
    footskip=1.2cm
]{geometry}

\usepackage{graphicx}
\usepackage{float}

\usepackage{xcolor}
\usepackage{colortbl}

\usepackage{array}
\usepackage{longtable}
\usepackage{tabularx}
\usepackage{booktabs}
\usepackage{multirow}

\usepackage{setspace}
\usepackage{enumitem}

\usepackage{titlesec}
\usepackage{titletoc}
\usepackage{tocloft}

\usepackage{fancyhdr}

\usepackage{caption}
\usepackage{subcaption}

\usepackage{amsmath}
\usepackage{amssymb}

\usepackage{url}
\usepackage{hyperref}
\usepackage{bookmark}

\usepackage{listings}
\usepackage{ragged2e}

%========================================================
% LANGUAGE SETTINGS
%========================================================

\setdefaultlanguage{english}
\setotherlanguage{arabic}

% English font
\setmainfont{TeX Gyre Termes}

% Arabic body text
\newfontfamily\arabicfont[
    Script=Arabic,
    Scale=1.08
]{Noto Naskh Arabic}

% Arabic title only
\newfontfamily\ardititlefont[
    Script=Arabic,
    Scale=1.05
]{Noto Kufi Arabic}

%========================================================
% ARDI GREEN THEME
%========================================================

\definecolor{ArdiGreen}{RGB}{0,102,51}
\definecolor{ArdiDarkGreen}{RGB}{0,82,41}
\definecolor{ArdiLightGreen}{RGB}{226,239,218}
\definecolor{ArdiVeryLightGreen}{RGB}{242,249,245}

\definecolor{HeaderGray}{HTML}{666666}
\definecolor{FooterGray}{HTML}{888888}
\definecolor{LightGray}{HTML}{CCCCCC}

%========================================================
% GLOBAL DOCUMENT SETTINGS
%========================================================

\onehalfspacing

\setlength{\parindent}{0pt}
\setlength{\parskip}{6pt}

\setlist[itemize]{
    leftmargin=1.5em,
    itemsep=3pt,
    topsep=5pt
}

\setlist[enumerate]{
    leftmargin=1.8em,
    itemsep=3pt,
    topsep=5pt
}

\justifying

%========================================================
% HYPERLINK SETTINGS
%========================================================

\hypersetup{
    colorlinks=true,
    linkcolor=ArdiGreen,
    citecolor=ArdiGreen,
    urlcolor=ArdiGreen,
    pdfauthor={
        Kinan Awartani,
        Taima Abu Bshara,
        Lujain Manasreh,
        Zaid Rajab,
        Amer Riad
    },
    pdftitle={
        Ardi - AI-Powered Real Estate Monitoring
        and Fraud Detection System
    }
}

%========================================================
% HEADER AND FOOTER
%========================================================

\pagestyle{fancy}

\fancyhf{}

\fancyhead[L]{
    \fontsize{9}{11}\selectfont
    \color{HeaderGray}
    \textbf{PTUK --- AI Graduation Project}
}

\fancyhead[R]{
    \fontsize{9}{11}\selectfont
    \color{ArdiGreen}
    \textit{Ardi}
}

\fancyfoot[L]{
    \fontsize{9}{11}\selectfont
    \color{FooterGray}
    Palestine Technical University -- Kadoorie
    \;|\;
    Faculty of IT \& AI
}

\fancyfoot[R]{
    \fontsize{9}{11}\selectfont
    \color{FooterGray}
    \thepage
}

\renewcommand{\headrulewidth}{0.4pt}
\renewcommand{\footrulewidth}{0.4pt}

\renewcommand{\headrule}{
    \hbox to\headwidth{
        \color{LightGray}
        \leaders\hrule height \headrulewidth\hfill
    }
}

\renewcommand{\footrule}{
    \hbox to\headwidth{
        \color{LightGray}
        \leaders\hrule height \footrulewidth\hfill
    }
}

%========================================================
% PLAIN PAGE STYLE
%========================================================

\fancypagestyle{plain}{
    \fancyhf{}

    \fancyhead[L]{
        \fontsize{9}{11}\selectfont
        \color{HeaderGray}
        \textbf{PTUK --- AI Graduation Project}
    }

    \fancyhead[R]{
        \fontsize{9}{11}\selectfont
        \color{ArdiGreen}
        \textit{Ardi}
    }

    \fancyfoot[L]{
        \fontsize{9}{11}\selectfont
        \color{FooterGray}
        Palestine Technical University -- Kadoorie
        \;|\;
        Faculty of IT \& AI
    }

    \fancyfoot[R]{
        \fontsize{9}{11}\selectfont
        \color{FooterGray}
        \thepage
    }

    \renewcommand{\headrulewidth}{0.4pt}
    \renewcommand{\footrulewidth}{0.4pt}
}

%========================================================
% CHAPTER STYLE
%========================================================

\titleformat{\chapter}[block]
    {\normalfont\fontsize{22}{26}\selectfont\bfseries\color{ArdiGreen}}
    {Chapter \thechapter:}
    {0.5em}
    {}

\titlespacing*{\chapter}
    {0pt}
    {-10pt}
    {22pt}

%========================================================
% SECTION STYLE
%========================================================

\titleformat{\section}
    {\normalfont\fontsize{16}{19}\selectfont\bfseries\color{ArdiGreen}}
    {\thesection}
    {0.8em}
    {}

\titlespacing*{\section}
    {0pt}
    {18pt}
    {8pt}

%========================================================
% SUBSECTION STYLE
%========================================================

\titleformat{\subsection}
    {\normalfont\fontsize{13}{16}\selectfont\bfseries\color{ArdiGreen}}
    {\thesubsection}
    {0.8em}
    {}

\titlespacing*{\subsection}
    {0pt}
    {14pt}
    {6pt}

%========================================================
% SUBSUBSECTION STYLE
%========================================================

\titleformat{\subsubsection}
    {\normalfont\fontsize{12}{15}\selectfont\bfseries\color{ArdiDarkGreen}}
    {\thesubsubsection}
    {0.8em}
    {}

%========================================================
% TABLE SETTINGS
%========================================================

\renewcommand{\arraystretch}{1.35}

\newcolumntype{Y}{
    >{\centering\arraybackslash}X
}

\newcolumntype{L}[1]{
    >{\raggedright\arraybackslash}p{#1}
}

\newcolumntype{C}[1]{
    >{\centering\arraybackslash}p{#1}
}

%========================================================
% FIGURE AND TABLE CAPTIONS
%========================================================

\captionsetup{
    font=small,
    labelfont=bf,
    justification=centering
}

\captionsetup[table]{
    position=bottom
}

%========================================================
% TABLE OF CONTENTS SETTINGS
%========================================================

\renewcommand{\contentsname}{Table of Contents}

\setlength{\cftbeforechapskip}{6pt}

\renewcommand{\cftchapfont}{
    \bfseries
}

\renewcommand{\cftchappagefont}{
    \bfseries
}

%========================================================
% CODE LISTING SETTINGS
%========================================================

\lstset{
    basicstyle=\ttfamily\small,
    breaklines=true,
    frame=single,
    numbers=left,
    numberstyle=\tiny,
    tabsize=4,
    showstringspaces=false
}

%========================================================
% PROJECT INFORMATION
%========================================================

\newcommand{\ProjectName}{Ardi}
\newcommand{\ProjectArabicName}{أرضي}

\newcommand{\ProjectFullTitle}{
AI-Powered Real Estate Monitoring and Fraud Detection System
}

\newcommand{\UniversityName}{
Palestine Technical University -- Kadoorie
}

\newcommand{\FacultyName}{
Faculty of Information Technology and Artificial Intelligence
}

\newcommand{\DepartmentName}{
Department of Artificial Intelligence
}

\newcommand{\BranchName}{
Tulkarm Branch
}

\newcommand{\SupervisorName}{
Dr. Areej Jaber
}

%========================================================
% BEGIN DOCUMENT
%========================================================

\begin{document}

%========================================================
% TITLE PAGE
%========================================================

\begin{titlepage}

\thispagestyle{empty}

\begin{center}

\vspace*{0.15cm}

%--------------------------------------------------------
% UNIVERSITY INFORMATION
%--------------------------------------------------------

{\fontsize{16}{20}\selectfont
\textbf{\UniversityName}
}

\vspace{0.15cm}

{\fontsize{14}{18}\selectfont
\textbf{\FacultyName}
}

\vspace{0.15cm}

{\fontsize{13}{17}\selectfont
\textbf{
\DepartmentName
\;|\;
\BranchName
}
}

\vspace{0.65cm}

%--------------------------------------------------------
% UNIVERSITY LOGO
%--------------------------------------------------------

\includegraphics[
    width=0.21\textwidth
]{ptuk_logo.png}

\vspace{0.7cm}

%--------------------------------------------------------
% PROJECT TITLE
%--------------------------------------------------------

{\fontsize{29}{35}\selectfont
\bfseries
\textcolor{ArdiGreen}{\ProjectName}
}

\vspace{0.08cm}

{\fontsize{25}{31}\selectfont
\bfseries
\textcolor{ArdiGreen}{
\textarabic{\ProjectArabicName}
}
}

\vspace{0.35cm}

{\fontsize{17}{21}\selectfont
\bfseries
\textcolor{ArdiDarkGreen}{
\ProjectFullTitle
}
}

\vspace{1cm}

%--------------------------------------------------------
% DEGREE STATEMENT
%--------------------------------------------------------

{\fontsize{12.5}{17}\selectfont
A Graduation Project Submitted in Partial Fulfillment of the
}

\vspace{0.10cm}

{\fontsize{12.5}{17}\selectfont
Requirements for the Bachelor's Degree in Artificial Intelligence
}

\vspace{0.9cm}

%--------------------------------------------------------
% PREPARED BY TABLE
%--------------------------------------------------------

\renewcommand{\arraystretch}{1.35}

\begin{tabular}{
|>{\centering\arraybackslash}p{7cm}|
>{\centering\arraybackslash}p{4cm}|
}

\hline

\rowcolor{ArdiGreen}

\textcolor{white}{
\textbf{Prepared By}
}
&
\textcolor{white}{
\textbf{Student ID}
}
\\

\hline

\rowcolor{ArdiLightGreen}
Kinan Awartani
&
202311151
\\

\hline

Taima Abu Bshara
&
202310481
\\

\hline

\rowcolor{ArdiLightGreen}
Lujain Manasreh
&
202310631
\\

\hline

Zaid Rajab
&
202310471
\\

\hline

\rowcolor{ArdiLightGreen}
Amer Riad
&
202220062
\\

\hline

\end{tabular}

\vspace{0.85cm}

%--------------------------------------------------------
% SUPERVISOR
%--------------------------------------------------------

{\fontsize{13}{17}\selectfont
\textbf{Supervised By}
}

\vspace{0.25cm}

{\fontsize{13}{17}\selectfont
\textbf{\SupervisorName}
}

\vspace{0.10cm}

{\fontsize{12}{16}\selectfont
Department of Artificial Intelligence
}

\vfill

%--------------------------------------------------------
% ACADEMIC YEAR
%--------------------------------------------------------

{\fontsize{12}{16}\selectfont
\textbf{
Academic Year 2026 / 2027 --- first Semester
}
}

\end{center}

\end{titlepage}

%========================================================
% FRONT MATTER
%========================================================

\pagenumbering{roman}

%========================================================
% DECLARATION OF ORIGINALITY
%========================================================

\chapter*{Declaration of Originality}

\addcontentsline{toc}{chapter}{Declaration of Originality}

We hereby declare that this graduation project, titled
``\ProjectName{} -- \ProjectFullTitle,''
is our own original work carried out under the supervision of
\SupervisorName{} at Palestine Technical University -- Kadoorie.

We further declare that:

\begin{enumerate}

    \item This work has not been submitted for any other degree
    at this or any other institution.

    \item All sources and references used have been properly
    cited and acknowledged.

    \item Any assistance received from others has been duly
    acknowledged.

    \item The use of generative AI tools, if any, has been
    disclosed and cited according to the department's
    academic integrity policy.

\end{enumerate}

\vspace{1cm}

\begin{center}

\renewcommand{\arraystretch}{1.7}

\begin{tabular}{
|L{6cm}|
C{3cm}|
C{3cm}|
}

\hline

\rowcolor{ArdiGreen}

\textcolor{white}{\textbf{Student Name}}
&
\textcolor{white}{\textbf{Signature}}
&
\textcolor{white}{\textbf{Date}}
\\

\hline

\rowcolor{ArdiVeryLightGreen}
Kinan Awartani
&
KINAN
&
7 october 2026
\\

\hline

Taima Abu Bshara
&
TAIMA
&
7 october 2026
\\

\hline

\rowcolor{ArdiVeryLightGreen}
Lujain Manasreh
&
LUJAIN
&
7 october 2026
\\

\hline

Zaid Rajab
&
ZAID
&
7 october 2026
\\

\hline

\rowcolor{ArdiVeryLightGreen}
Amer Riad
&
AMER
&
7 october 2026
\\

\hline

\end{tabular}

\end{center}

\clearpage

%========================================================
% ABSTRACT
%========================================================

\chapter*{Abstract}

\addcontentsline{toc}{chapter}{Abstract}

\textbf{Title:}
\ProjectName{} -- \ProjectFullTitle

\vspace{0.3cm}

\textbf{Keywords:}
Artificial Intelligence,
Real Estate,
Fraud Detection,
Anomaly Detection,
Machine Learning,
Data Analytics

\vspace{0.5cm}

Ardi is an AI-powered real estate monitoring and fraud detection
system designed to support authorities in monitoring and analysing
property and land transactions.

The system centralises real estate information including properties,
buyers, sellers, transaction values, geographical locations,
historical records, and transaction dates.

Using data analytics and machine learning techniques, the system
analyses historical and current transaction behaviour to detect
unusual patterns and anomalies that may require further investigation.

Rather than automatically classifying individuals or transactions as
fraudulent, Ardi operates as a decision-support system that assigns
an anomaly or suspicion score to unusual activities and presents the
results to authorised users for human review.

The platform also provides an interactive dashboard, geographical
visualisation of real estate activity, and periodic reports to assist
authorities in understanding market behaviour and identifying
transactions that deserve further attention.

\clearpage

%========================================================
% ARABIC ABSTRACT
%========================================================

\chapter*{\textarabic{الملخص}}

\addcontentsline{toc}{chapter}{Arabic Abstract}

\begin{Arabic}

\textbf{عنوان المشروع :}

\textbf{أرضي هو نظام ذكي لمراقبة العقارات واكتشاف الأنماط المحتملة للاحتيال}

\vspace{0.3cm}

\textbf{الكلمات المفتاحية:}

الذكاء الاصطناعي،
العقارات،
كشف الاحتيال،
اكتشاف الشذوذ،
التعلم الآلي،
تحليل البيانات

\vspace{0.5cm}

مشروع أرضي هو نظام ذكي قائم على تقنيات الذكاء الاصطناعي
يهدف إلى مساعدة الجهات المختصة في مراقبة وتحليل معاملات الأراضي
والعقارات.

يقوم النظام بجمع وتنظيم بيانات العقارات والمعاملات والأطراف
المرتبطة بها والأسعار والتواريخ والمواقع الجغرافية ضمن منصة
مركزية، ثم يستخدم تقنيات تحليل البيانات والتعلم الآلي لتحليل
النشاط العقاري الحالي والتاريخي واكتشاف الأنماط والسلوكيات
غير الطبيعية التي قد تستحق المزيد من المراجعة والتحقيق.

لا يقوم النظام بإصدار حكم قطعي بأن معاملة معينة احتيالية،
بل يعمل كنظام لدعم القرار من خلال تحديد الحالات غير الطبيعية
وإعطائها درجة اشتباه أو درجة شذوذ تساعد الموظف المختص على
ترتيب أولويات المراجعة.

كما يوفر النظام لوحة تحكم تفاعلية، وتحليلاً جغرافياً للنشاط
العقاري، وتقارير دورية تهدف إلى تحويل البيانات الخام إلى
معلومات ومؤشرات تساعد الجهات المختصة في عملية المراقبة
واتخاذ القرار.

\end{Arabic}

\clearpage

%========================================================
% DEDICATION
%========================================================

\chapter*{Dedication}

\addcontentsline{toc}{chapter}{Dedication}

\vspace{2cm}

\begin{center}

\itshape

To our families, whose endless support and patience made this work
possible.

\vspace{0.7cm}

To our supervisor and professors at Palestine Technical University --
Kadoorie who guided us throughout this journey.

\vspace{0.7cm}

And to Palestine.

\end{center}

\clearpage

%========================================================
% ACKNOWLEDGEMENTS
%========================================================

\chapter*{Acknowledgements}

\addcontentsline{toc}{chapter}{Acknowledgements}

We would like to express our sincere gratitude to
\SupervisorName{} for her invaluable guidance, continuous support,
constructive feedback, and encouragement throughout the development
of this graduation project.

We are deeply grateful to the Faculty of Information Technology and
Artificial Intelligence at Palestine Technical University -- Kadoorie
for providing the academic environment and resources necessary to
complete this work.

We also extend our sincere appreciation to all instructors,
colleagues, friends, and individuals who contributed directly or
indirectly to this project.

Finally, we would like to thank our families for their continuous
support, patience, and encouragement throughout our academic journey.

\clearpage

%========================================================
% TABLE OF CONTENTS
%========================================================

\tableofcontents

\clearpage

%========================================================
% LIST OF FIGURES
%========================================================

\listoffigures

\addcontentsline{toc}{chapter}{List of Figures}

\clearpage

%========================================================
% LIST OF TABLES
%========================================================

\listoftables

\addcontentsline{toc}{chapter}{List of Tables}

\clearpage

%========================================================
% MAIN MATTER
%========================================================

\pagenumbering{arabic}

\setcounter{page}{1}

%========================================================
% CHAPTER 1
%========================================================

\chapter{Introduction}

This chapter introduces the research problem, provides background
context, and outlines the objectives, research questions, significance,
scope, limitations, and overall structure of the report.

%--------------------------------------------------------
\section{Background and Motivation}
%--------------------------------------------------------

Real estate and land transactions generate large amounts of data
related to properties, buyers, sellers, transaction values,
geographical locations, and historical ownership records.

As the number of transactions increases, manually reviewing and
monitoring this information becomes more difficult, especially when
data is distributed across multiple sources and systems.

Artificial intelligence and data analytics provide opportunities to
analyse historical and current real estate activity automatically,
identify recurring patterns, and highlight unusual behaviour that may
require additional investigation.

The Ardi project is motivated by the need for a centralised intelligent
platform that supports authorities in monitoring property transactions
more efficiently and converting large amounts of raw transaction data
into useful information.

%--------------------------------------------------------
\section{Problem Statement}
%--------------------------------------------------------

The increasing volume and complexity of real estate transactions make
manual monitoring time-consuming and inefficient.

A transaction may appear normal when examined independently, while a
historical, relational, or geographical analysis may reveal unusual
behaviour.

Examples include significant changes in property prices, frequent
transactions between the same parties, repeated resale of the same
property within short periods, unusually high transaction activity by
a specific person or organisation, or significant differences between
the prices of similar properties.

Identifying such patterns manually requires substantial effort and may
be difficult when large volumes of data are involved.

Therefore, there is a need for a decision-support system that can
centralise real estate data, analyse transaction behaviour, identify
unusual patterns, and present suspicious cases to authorised users for
further review.

%--------------------------------------------------------
\section{Project Objectives}
%--------------------------------------------------------

The main objectives of the Ardi project are:

\begin{enumerate}

    \item To develop a centralised platform for managing property
    and real estate transaction data.

    \item To analyse historical and current real estate activity.

    \item To apply artificial intelligence and machine learning
    techniques for anomaly detection.

    \item To identify transactions and behavioural patterns that
    may require further review.

    \item To analyse real estate activity geographically.

    \item To provide an interactive dashboard containing useful
    statistics and key performance indicators.

    \item To generate daily, weekly, monthly, and yearly reports.

    \item To reduce the time and effort required for manual
    transaction monitoring.

    \item To transform raw real estate data into useful information
    that supports decision-making.

\end{enumerate}

%--------------------------------------------------------
\section{Research Questions}
%--------------------------------------------------------

This project aims to answer the following research questions:

\begin{enumerate}

    \item How can artificial intelligence techniques be used to
    detect unusual patterns in real estate transactions?

    \item Which transaction, property, temporal, geographical, and
    behavioural features are most useful for identifying anomalous
    real estate activities?

    \item How can anomaly detection results be presented in a way
    that supports human review and decision-making?

    \item How can geographical and historical transaction analysis
    improve the monitoring of real estate activities?

\end{enumerate}

%--------------------------------------------------------
\section{Significance of the Study}
%--------------------------------------------------------

The significance of this project lies in integrating real estate data,
artificial intelligence, data analytics, geographical analysis,
interactive visualisation, and automated reporting within one
decision-support platform.

Instead of relying entirely on manual inspection, authorised users can
use Ardi to identify transactions and activities that deserve greater
attention.

The system can also provide a clearer understanding of market trends,
transaction activity, unusual behaviour, and geographical
concentrations.

%--------------------------------------------------------
\section{Scope and Limitations}
%--------------------------------------------------------

The scope of Ardi includes property and transaction data management,
historical transaction analysis, anomaly detection, dashboard
visualisation, geographical analysis, and periodic reporting.

The system is intended to support authorised authorities and
specialised employees rather than general public users.

The system does not automatically determine that a person or
transaction is fraudulent.

It does not replace human investigation, legal procedures, or official
decision-making.

The quality and accuracy of the system depend on the availability,
quality, completeness, and accuracy of the provided data.

The performance of the machine learning model may also be affected by
limited, imbalanced, or incomplete training data.

%--------------------------------------------------------
\section{Report Organisation}
%--------------------------------------------------------

The remainder of this report is organised as follows:

\begin{itemize}

    \item Chapter 2 presents the theoretical background, related
    work, and research gap.

    \item Chapter 3 describes the methodology, system requirements,
    architecture, database design, data processing, and AI approach.

    \item Chapter 4 presents the system implementation and
    experimental results.

    \item Chapter 5 discusses the results, strengths, limitations,
    and challenges.

    \item Chapter 6 concludes the project and presents possible
    future improvements.

\end{itemize}

%========================================================
% END DOCUMENT
%========================================================

\end{document}